\bs
\ni{\bf Lemma 4.1.20} If $c\in V_k$ and $d\in V_{k'}$, then
\begin{verse}
(i) \ $\a_{d,j+k}\a_{c,i}=\a_{c,i}\a_{d,j}$ \ for \ 
$i\leq j$.

(ii) For \ $i < j$, \ $a_{j-1}a_i=a_ia_j$,  \
 $\mu_{j-1}\mu_i=\mu_i\mu_j$. \newline
For \ $i\leq j$, \ $\eta_{j+1}\eta_i=\eta_i\eta_j$, \
$\kappa_{j+1}\kappa_i=\kappa_i\kappa_j$.

(iii) For $i\leq j$,
\begin{eqnarray*}
   a_{j+k}\a_{c,i}=\a_{c,i}a_j &
a_{j-1}\a_{c,j}=\a_{c,j-1}a_{i-1} \\
\mu_{j+k+1}\a_{c,i}=\a_{c,i}\mu_{j+1} &
\mu_{i-1}\a_{c,j}=\a_{c,j-1}\mu_{i-1} \\
\eta_{j+k}\a_{c,i}=\a_{c,i}\eta_j &
\eta_{i-1}\a_{c,j-1}=\a_{c,j}\eta_{i-1} \\
\kappa_{j+k}\a_{c,i}=\a_{c,i}\kappa_j & \kappa_{i-1}\a_{c,j}=\a_{c,j+1}
\kappa_{i-1}
\end{eqnarray*}

(iv) For $i < j$,
\begin{eqnarray*}
\mu_j a_i =a_i\mu_{j+1} & \mu_ia_j=a_{j-1}\mu_i\\
\eta_{j-1}a_i=a_i\eta_j &   \eta_{i+1}a_j=a_{j+1}\eta_{i+1}\\
\kappa_{j-1}a_i=a_i\kappa_j &  \kappa_ia_j=a_{j+1}\kappa_i\\
\eta_{j-1}\mu_i=\mu_i\eta_j &  \eta_i\mu_j=\mu_{j+1}\eta_i\\
\kappa_{j-2}\mu_i=\mu_i\kappa_{j-1} &  \kappa_i\mu_j=\mu_{j+1}\kappa_i\\
\kappa_j\eta_i =\eta_i\kappa_{j-1} &  \kappa_i\eta_j=\eta_{j+1}\kappa_i
\end{eqnarray*}

(v) 
\begin{eqnarray*}
a_i\kappa_i={\op{id}} \\
\mu_i\eta_i={\op{id}} \\
a_{i+1}\kappa_i={\op{id}} \\
\mu_{i+1}\eta_i={\op{id}} 
\end{eqnarray*}

(vi) 
\begin{eqnarray*}
a_j\eta_j=\d{\op{id}} \\
\mu_{j+1}\kappa_j=\d{\op{id}}.
\end{eqnarray*}
\end{verse}
(These identities hold when $i, \ j$ and $k$ are such that all the
maps involved are defined by 4.1.19.)

\bigskip
{\sc Proof.} Almost all cases of identities (i)--(iv) are trivial
as they  can be written so as to involve distant tensor product
indices: thus they just amount to a renumbering.
The ones that involve some interaction between the tensor product
components are
$$
\mu_i\mu_i=\mu_i\mu_{i+1} \ , \ \
\kappa_{i+1}\kappa_i=\kappa_i\kappa_i \ , \ \
\kappa_{i-1}\mu_i=\mu_i\kappa_i \ , \ \
\kappa_i\mu_{i+1}=\mu_{i+2}\kappa_i \ .
$$
These all follow easily from associativity of multiplication and\\
$\dsize{\sum_{b\in B}\!b\!\otimes\! b^*x\!=\!\sum_{b\in B}}xb \otimes b^*$
for $x\in M$, which is 4.1.2.

\begin{eqnarray*}
{\mbox{For (v):}}\ a_i\kappa_i& = &{\op{id}}\ {\mbox{follows from}} \sum_{b\in B} E_N(xb)b^*=x \ \ \ (4.1.4)\\
\mu_i\eta_i & = & {\op{id}}\  {\mbox{follows from}}  x1=x\\
a_{i+1}\kappa_i & = & {\op{id}}\  {\mbox{follows from}} 
\sum_{b\in B} bE(b^*x)=x \ \ \ \ (4.1.4)\\
\mu_{i+1}\eta_i & = &{\op{id}}\   {\mbox{follows from}}\  x1=x
\end{eqnarray*}

\begin{eqnarray*}
{\mbox{For (vi):}}\   a_j\eta_j & = & \d{\op{id}}\   {\mbox{follows from}}\ E_N(1)=1\\
\mu_{j+1}\kappa_j&=&\d{\op{id}}\  {\mbox{follows from}}\ \sum_{b\in B} bb^*=\d^2{\op{id}} \ \ \ (4.1.3)
\end{eqnarray*}
\qed

\medskip
\ni{\bf Lemma 4.1.21} \ If \ $x\in M$, \
$2\leq r\leq k$, then \
$v_kxv_r=v_{r-2}xv_k$ (where $v_0=1$), and
$v_kxv_1=\d E_N(x)v_k$.

\bigskip
{\sc Proof.} Simple manipulation of 4.1.8 and 4.1.9. \qed

\medskip
\ni{\bf Lemma 4.1.22} \ If \ $x\in \otimes^k_N M$, \ then \
\begin{verse}
(i)  $\th^{-1}E_{M_{k-2}}\th(x)=\d^{-1}
a_{m+1}(x)$ \ if \ $k$ is odd, $k=2m+1$.

(ii) $\th^{-1}E_{M_{k-2}}\th(x)=\d^{-1}\mu_{m+1}(x)$ \ if \ $k$ is even, $k=2m$.
\end{verse}

\bigskip
{\sc Proof.} Let $x$ be of the form $x_1\otimes
x_2\otimes\dots\otimes x_k$. Then
$$
E_{M_{k-2}}(\th(x))=
E_{M_{k-2}}(x_1v_1x_2v_2\dots v_{k-1}x_k) =
\d^{-1}x_1v_1x_2v_2\dots x_{k-2}v_{k-2}x_{k-1}v_{k-2}x_k
$$

\ni${\underline{\mbox{Case (i)}}}$. If \ $k=2m+1$ we may apply
4.1.21
$m\!-\!1$ times to obtain
\begin{eqnarray*}
E_{M_{k-2}}(\th(x)) & = \d^{-1}x_1v_1\dots x_mv_m x_{m+1}v_1
x_{m+2}v_{m+1}\dots x_{k-1}v_{k-2}x_k\\
&= x_1v_1\dots x_mv_m E(x_{m+1})x_{m+2}v_{m+1}\dots
x_{k-1}v_{k-2}x_k \\
&= \d^{-1}\th(a_{m+1}(x))
\end{eqnarray*}
${\underline{\mbox{Case (ii)}}}$. If \ $k=2m$ we  apply 4.1.21
$m\!-\!1$ times to obtain
\begin{eqnarray*}
E_{M_{k-2}}(\th(x)) & = \d^{-1}x_1v_1\dots x_mv_m x_{m+1}v_0
 x_{m+2}v_{m+1}\dots x_{k-1}v_{k-2}x_k\\
&= \d^{-1}\th(\mu_{m+2}(x)) \quad {\text{(since}} \ v_0\!=\!1).
\end{eqnarray*}
\qed

\bs
\ni{\bf Lemma 4.1.23} \ If \ $y\in \otimes^k_N M$ \ and \
$x,z\in M$, then
$$
xv^*_k\th(y)v_{k+1}z=\th(x\otimes y\otimes z)
$$

\bs
{\sc Proof.} Simple commutation of $E_i$ with $M$ and $E_j$'s
$j\leq i-2$. \qed

\bs
In the next lemma, let $\chi_j=a_j\mu_{j+1}$, \
$j=1,2,\dots k-1$.

\bs
\ni{\bf Lemma 4.1.24} \ If \ $x,y\in \otimes^k_N M$, \ then \
$\th(x)\th(y)=$
\begin{verse}
(i) $\th(\chi_{m+1}\chi_{m+2}\dots\chi_k(x\otimes_N y))$ \
if $k$ is even, $k=2m$.

(ii) $\th(\mu_{m+2}\chi_{m+2}\dots\chi_k(x\otimes_N y))$ \
if $k$ is odd, $k=2m+1$.
\end{verse}

\bs
{\sc Proof.}  By induction on $k$. Let $x=x_1\otimes
x_2\dots\otimes x_{k+1}$, \ $y=y_1\otimes
y_2\dots\otimes y_{k+1}$, \ then

\begin{eqnarray*}
\th(x)\th(y) & = & x_1v^*_k kx_2v^*_{k-1}\dots v^*_1x_{k+1}y_1v_1y_2v_2 \dots y_kv_ky_{k+1} \\
&= & \d x_1v^*_{k-1} E_{M_{k-2}}(\th(x_2\otimes x_3\otimes\dots\otimes x_{k+1})\th(y_1\otimes\dots\otimes y_k))
v_ky_{k+1}
\end{eqnarray*}

(i) If $k$ is even, $k=2m$, by the inductive hypothesis we have
\begin{eqnarray*}
\th(x)\th(y) & = &\d x_1v^*_{k-1} E_{M_{k-2}} (\th(\chi_{m+1}\chi_{m+2}\dots\chi_k(x_2\otimes\dots\otimes x_{k+1}\otimes y_1\otimes\dots\otimes y_k)))v_ky_{k+1} \\
& = & x_1v^*_{k-1} \th(\mu_{m+1}\chi_{m+1}\dots\chi_k (x_2\otimes\dots\otimes y_k))v_ky_{k+1} \ \ \mbox{(by\  Lemma\ 4.1.22)}\\
& = &\th(x_1\otimes (\mu_{m+1}\chi_{m+1}\chi_{m+2}\dots \chi_k(x_2\otimes\dots\otimes y_k)\otimes y_{k+1})) \ \ \mbox{(by Lemma 4.1.23)}\\
& = &\th(\mu_{m+2}\chi_{m+2}\chi_{m+3}\dots \chi_{k+1}(x_1\otimes x_2\otimes x_{k+1}\otimes y_1 \otimes\dots\otimes y_{k+1}))
\end{eqnarray*}

(ii) If $k$ is odd, $k=2m+1$,
\begin{eqnarray*}
\th(x)\th(y) & = & x_1 v^*_{k-1} E_{M_{k-2}} (\th(\mu_{m+2}\chi_{m+2}\chi_{m+3}\dots\chi_k (x_2\otimes\dots x_{k+1}\otimes y_1 \otimes\dots\otimes y_k)))v_ky_{k+1} \\
& & \mbox{(by the induction hypothesis)} \\
& = & x_1v^*_{k-1}\th(a_{m+1}\mu_{m+2}\chi_{m+2}\chi_{m+3}\dots \chi_k(x_2\otimes\dots\otimes y_k))v_ky_{k+1}\\
& & \mbox{(by Lemma 4.1.22)}\\
& = & \th(x_1\otimes \chi_{m+1}\chi_{m+2}\dots \chi_k(x_2\otimes\dots\otimes y_k)\otimes y_{k+1})) \quad \mbox{(by Lemma 4.1.23)}\\
& = &\th(\chi_{m+2}\chi_{m+3}\dots \chi_{k+1}(x_1\otimes x_2 \dots\otimes y_{k+1})) \ .
\end{eqnarray*}
It only remains to check the formula for $k\!=\!1, \ (m\!=\!0)$.
Then $\th(x)=x$, \ $\th(y)=y$ and the formula reads
\[
\th(x)\th(y)=xy=\th(\mu_2(x\otimes y))=\th(xy) \ .
\]
\qed

\bs
\ni{\bf Lemma 4.1.25}
For $m=1,2,\dots$, and $x_{m+1},x_{m+2},\dots x_{2m}\in M$,
$$
v_mv_{m+1}x_{m+1}v_{m+2}x_{m+2}v_{m+3}\dots
v_{2m}x_{2m}=v_mx_{m+1}v_{m+1}x_{m+2}\dots v_{2m-1}x_{2m} \  .
$$

\bigskip
{\sc Proof.} Induction on $m$.

For $m=1$ the formula reads $E_1E_2E_1x_2=E_1x_2$ which is correct.
Now suppose the formula holds for $m$, then
$$
v_{m+1}v_{m+2}x_{m+2} =v_{m+1}x_{m+2}E_{m+2}E_{m+1}\dots E_3 \qquad
\mbox{(by 4.1.8)}
$$
so
$$
v_{m+1}v_{m+2}x_{m+2} v_{m+3}x_{m+3}\dots v_{2m+2}x_{2m+2} =
(v_{m+1}x_{m+2}) V_mV_{m+1} y_{m+1}V_{m+2}y_{m+2}\dots
V_{2m}y_{2m}
$$
where \  $V_n=E_{n+2}E_{n+1}\dots E_3$ \ and \
$y_n=E_2E_1x_{n+2}$.

We may now apply the inductive hypothesis to the subfactor
$M_1\subset M_2$ (for which the $E_i$'s are just those for
$N\subset M$, shifted by 2), to obtain
$$
v_{m+1}x_{m+2}V_my_{m+1}\dots V_{2m}y_{2m}=
v_{m+1}x_{m+2}v_{m+2}x_{m+3}\dots v_{2m+1}x_{2m+2}
$$
\qed

\bs
\ni{\bf Corollary 4.1.26} \ With notation as above, for \
$m=2,3,4,\dots,$
\begin{eqnarray*}
& E_mE_{m-1}\dots E_2(v_{m+1}x_{m+1})(v_{m+2}x_{m+2})\dots
(v_{2m-1}x_{2m-1}) \\
&\quad =
(v_m\dots x_{m+1})(v_{m+1}x_{m+2})\dots (v_{2m-2}x_{2m-1})
\end{eqnarray*}

\bs
{\sc Proof.} Write \ $V_n=E_{n+1}E_{n-1}\dots E_2$, \
$y_n=E_1x_{n+1}$ and apply the 4.1.25 to the subfactor $M\subset
M_1$. \qed

\bs
\ni{\bf Lemma 4.1.27} \ For
$p=1,2,\dots,$
\[
\sum_{b_1,b_2,\dots b_p\in B} (b_1v_1)(b_2v_2)\dots
(b_pv_p)v_{p+1}v_{p+2} b^*_{p}v_{p+3}b^*_{p-1} v_{p+4}\dots
v_{2p+1}b^*_1 = \d^p E_{2p+1}
\]

\bs
{\sc Proof.}
By induction on $p$. For $p=1$ the formula reads
$$
\sum_{b\in B} bE_1E_2E_1E_3E_2E_1b^* = E_3 \sum_{b\in B}
bE_1b^* = \d E_3  \ ,
$$
which is correct. Now observe that
$$
\sum_b bv_p v_{p+1}v_{p+2}b^*v_{p+3}=\d
E_p E_{p-1}\dots E_2 V_{p-1}V_p V_{p+1} E_2E_1
$$
where $V_n=E_{n+2}E_{n+1}\dots E_3$, and that
$$
b_1v_1b_2v_2\dots b_{p-1}v_{p-1}E_pE_{p-1}\dots E_2=
y_1V_1y_2V_2\dots y_{p-2}V_{p-2}y_{p-1}
$$
where $y_i=b_iE_1E_2$, so that
\begin{eqnarray*}
&\dsize{\sum_{b_1,\dots ,b_p\in B} b_1v_1b_2v_2\dots b_pv_p v_{p+1}v_{p+2}b^*_p v_{p+3}b^*_{p-1} v_{p+4}\dots v_{2p+1}b^*_1} \\
& = \dsize{\d \sum_{y_1\dots y_{p-1}\in BE_1E_2} y_1V_1y_2V_2\dots y_{p-2}V_{p-2}y_{p-1}V_{p-1}V_pV_{p+1} y^*_{p-1}V_{p+2}\dots V_{2p-1}y^*_1}
\end{eqnarray*}
so since $\{ bE_1E_2\mid b\in B\}$ is a basis for $M_2$ over $M_1$,
we are through by induction. \qed

\ni{\bf Corollary 4.1.28} \ For
$p=1,2,\dots,$
$$
\sum_{b_1,b_2,\dots ,b_{p+1}\in B}
(b_1v_1)(b_2v_2)\dots (b_pv_p)b^*_p b_{p+1}v_{p+1}
b^*_{p+1}v_{p+2} b^*_{p-1}\dots v_{2p}b^*_1 =\d^{p+1}E_{2p}
$$

\bigskip
{\sc Proof.} Observe that
$$
\sum_{b_p,b_{p+1}\in B} b_pv_p b^*_pb_{p+1} v_{p+1}
b^*_{p+1}v_{p+2} = \d^2 V_{p-1}V_p V_{p+1}E_1 \ ,
$$
where $V_n=E_{n+1}E_n\dots E_2$, so the left-hand side of the
equation becomes
$$
\d^2\sum_{y_1,\dots ,y_{p-1}} y_1V_1y_2V_2\dots y_{p-1}
V_{p-1} V_pV_{p+1} y^*_{p-1}V_{p+2}\dots V_{2p-1}y^*_1
$$
where $y_n=b_nE_1$, and this is $\d^{p+1}E_{2p}$ by 4.1.27 applied
to the subfactor $M\subset M_1$ with basis $\{bE_1\mid b\in B\}$.
\qed

\bs\bs\noindent{\bf 4.2 \ Subfactors give planar algebras}

\smallskip We keep the notation of $\S$4.1.
The next theorem legitimizes the use of pictures to prove
subfactor results.

\ni{\bf Theorem 4.2.1} Let $N\subset M$ be an extremal type
$II_1$ subfactor with $[M:N]^{\frac 12}=\d <\i$. For each $k$ let
$P^{N\subset M}_k=N'\cap M_{k-1}$ (isomorphic via $\th$ to $V_k$,
i.e.
$N$-central vectors in $\bigotimes^k_N M$). Then
$P^{N\subset M}=\bigcup_k P^{N\subset M}_k$ has a
spherical
$C^*$-planar algebra structure (with labelling set $P^{N\subset
M}$) for which
$\Phi(\tangle{x})=x$ and, suppressing the
presenting map~$\Phi$,
\begin{verse}
(i) For $i=1,2,\dots k-1,\boxj = E_i$ \bs

(ii) $\boxh{x} = \d E_{M'}(x)$ , \ \
 $\boxl{x} \ = \d E_{M_{k-2}}(x)$ \bs

(iii) $\tangle{x}\ \  \vln = \tangle{x}$ 
(where on the right, $x$ is considered as an element of $M_{k+1}$)
\bs

(iv) $Z(\boxrloop{x})=\d^k {\op{tr}}(x)$ \ \
$(x\in P^{N\subset M}_k)$
\end{verse}
\bs
Moreover, any other spherical planar algebra structure $\Phi'$ with
$\Phi'(\tangle{x})=x$ and {\rm (i),(ii),(iv)} for
$\Phi'$ is equal to $\Phi$.

\bigskip
{\sc Proof.}  The idea of the proof is fairly simple but it will
involve a lot of details, so we begin with an informal description
of the idea. We must show how to associate an element $\Phi(T)$ in
$N'\cap M_{k-1}$ to a tangle $T$  whose boxes are labelled by
elements of the appropriate $N'\cap M_j$. An example with $k=5$ is
given in Figure 4.2.2.
%%%%%%%%%%%%%%%%%%%%%%%%%%%%%%%%%%%%%%%%%%%%%%%%%%%%%%%%%%%%%%%%%%%
\[
	\begin{picture}(0,0)%
\epsfig{file=ps/pic84.ps}%
\end{picture}%
\setlength{\unitlength}{0.00041700in}%
%
\begingroup\makeatletter\ifx\SetFigFont\undefined
% extract first six characters in \fmtname
\def\x#1#2#3#4#5#6#7\relax{\def\x{#1#2#3#4#5#6}}%
\expandafter\x\fmtname xxxxxx\relax \def\y{splain}%
\ifx\x\y   % LaTeX or SliTeX?
\gdef\SetFigFont#1#2#3{%
  \ifnum #1<17\tiny\else \ifnum #1<20\small\else
  \ifnum #1<24\normalsize\else \ifnum #1<29\large\else
  \ifnum #1<34\Large\else \ifnum #1<41\LARGE\else
     \huge\fi\fi\fi\fi\fi\fi
  \csname #3\endcsname}%
\else
\gdef\SetFigFont#1#2#3{\begingroup
  \count@#1\relax \ifnum 25<\count@\count@25\fi
  \def\x{\endgroup\@setsize\SetFigFont{#2pt}}%
  \expandafter\x
    \csname \romannumeral\the\count@ pt\expandafter\endcsname
    \csname @\romannumeral\the\count@ pt\endcsname
  \csname #3\endcsname}%
\fi
\fi\endgroup
\begin{picture}(5144,3344)(1779,-3983)
\put(4051,-3286){\makebox(0,0)[lb]{\smash{\special{ps:gsave currentpoint currentpoint translate
-90.0 rotate neg exch neg exch translate}\SetFigFont{12}{14.4}{rm}$R_3$\special{ps:currentpoint grestore moveto}}}}
\put(4276,-2386){\makebox(0,0)[lb]{\smash{\SetFigFont{12}{14.4}{rm}$R_2$}}}
\put(3076,-2386){\makebox(0,0)[lb]{\smash{\SetFigFont{12}{14.4}{rm}$R_1$}}}
\end{picture}

\]
\begin{center}
Figure 4.2.2
\end{center}
%%%%%%%%%%%%%%%%%%%%%%%%%%%%%%%%%%%%%%%%%%%%%%%%%%%%%%%%%%%%%%%%%%%
\bs Shade the regions black and white and observe that a smooth
oriented curve starting and ending on the left-hand boundary, and
missing the internal boxes, will generically pass through a
certain number of black regions. Number the connected components of
the intersection of the curve with the black
 regions $1,2,\dots
n$ in the order they are crossed.
The regions on the curve will be used to index the tensor product
components in
$\bigotimes^n_N M$.

We will start with our curve close to the boundary so that it
crosses no black regions, and allow it to bubble outwards until it
is very close to the outside boundary at which point it will cross
$k$ black regions. As the curve bubbles out, it will pass through
non-generic  situations with respect to the strings of the tangle,
and it will envelop internal boxes. At generic times we will
associate an $N$-central element of $\bigotimes^n_N M$ with the
curve. As the curve passes through exceptional situations we will
change the element of $\bigotimes^n_N M$ according to certain
rules, the main one of which being that, when the curve envelops a
box labelled by a tensor, we will insert that label into the
tensor on the curve at the appropriate spot, as illustrated in
Figure~\ref{pic85}
%%%%%%%%%%%%%%%%%%%%%%%%%%%%%%%%%%%%%%%%%%%%%%%%%%%%%%%%%%%%%%%%%%%
\[
        \begin{picture}(0,0)%
\epsfig{file=ps/pic85.ps}%
\end{picture}%
\setlength{\unitlength}{0.00041700in}%
%
\begingroup\makeatletter\ifx\SetFigFont\undefined
% extract first six characters in \fmtname
\def\x#1#2#3#4#5#6#7\relax{\def\x{#1#2#3#4#5#6}}%
\expandafter\x\fmtname xxxxxx\relax \def\y{splain}%
\ifx\x\y   % LaTeX or SliTeX?
\gdef\SetFigFont#1#2#3{%
  \ifnum #1<17\tiny\else \ifnum #1<20\small\else
  \ifnum #1<24\normalsize\else \ifnum #1<29\large\else
  \ifnum #1<34\Large\else \ifnum #1<41\LARGE\else
     \huge\fi\fi\fi\fi\fi\fi
  \csname #3\endcsname}%
\else
\gdef\SetFigFont#1#2#3{\begingroup
  \count@#1\relax \ifnum 25<\count@\count@25\fi
  \def\x{\endgroup\@setsize\SetFigFont{#2pt}}%
  \expandafter\x
    \csname \romannumeral\the\count@ pt\expandafter\endcsname
    \csname @\romannumeral\the\count@ pt\endcsname
  \csname #3\endcsname}%
\fi
\fi\endgroup
\begin{picture}(11519,4162)(1704,-3976)
\put(10652,-1862){\makebox(0,0)[lb]{\smash{\special{ps:gsave currentpoint currentpoint translate
-30.0 rotate neg exch neg exch translate}\SetFigFont{12}{14.4}{rm}$x\otimes y\otimes z$\special{ps:currentpoint grestore moveto}}}}
\put(9452,-962){\makebox(0,0)[lb]{\smash{\SetFigFont{12}{14.4}{rm}$a$}}}
\put(9902,-3287){\makebox(0,0)[lb]{\smash{\SetFigFont{12}{14.4}{rm}$b$}}}
\put(4276,-1861){\makebox(0,0)[lb]{\smash{\special{ps:gsave currentpoint currentpoint translate
-30.0 rotate neg exch neg exch translate}\SetFigFont{12}{14.4}{rm}$x\otimes y\otimes z$\special{ps:currentpoint grestore moveto}}}}
\put(3076,-961){\makebox(0,0)[lb]{\smash{\SetFigFont{12}{14.4}{rm}$a$}}}
\put(3526,-3286){\makebox(0,0)[lb]{\smash{\SetFigFont{12}{14.4}{rm}$b$}}}
\end{picture}

\]
\begin{center}
Figure 4.2.3: Tensor on bubbling curve: $\dots\otimes           
a\otimes b\otimes\dots
\ \to \ \dots\otimes a\otimes x \otimes
y\otimes z \otimes b\otimes\dots$
\end{center}
%%%%%%%%%%%%%%%%%%%%%%%%%%%%%%%%%%%%%%%%%%%%%%%%%%%%%%%%%%%%%%%%%%

\bs\ni When the curve arrives very close to the boundary it will
have associated to it a central element of $\bigotimes^k_N M$ which
gives an element of $N'\cap M_{k-1}$ via $\th$. This element will
be $\Phi(T)$.

This strategy meets several obstacles.
\begin{verse}
(1) We must show that $\Phi(T)$ is well defined -- note that the
insertions of 4.2.3 are {\bf not} well defined for the tensor
product over $N$.

(2) $\Phi(T)$ must be central. This will require either
enveloping the boxes only starting from the white region touching
the first boundary point, or projecting onto central vectors at
each step. We will adopt the former policy.

(3)$\Phi(T)$ must be independent of isotopy of $T$ and the
choice of the path. Since we must show isotopy invariance
eventually, we might as well suppose that the tangles are in a
convenient standard form, since any two ways of arriving at that
standard form from a given $T$ will only differ by an isotopy.
\end{verse}

We begin the formal proof by describing the standard form.

\medskip A $k${\it -picture} (or just ``picture"
 if the value of $k$ is clear) will be the intersection of the unit
square $[0,1]\x [0,1]$ in the $x\!-\!y$ plane with a system of
smooth curves, called strings, meeting only in finitely many
isolated  singularities, called ``cusps", where $2m$ strings meet
in a single point. The strings must meet the boundary of $[0,1]\x
[0,1]$ transversally in just $2k$ points on the boundary line
$[0,1]\x
\{1\}$.  A cusp $(x,y)$ will be said to be in {\it standard form} \
if, in some neighborhood of $(x,y)$ the $y$-coordinates of points
on the strings are all greater than $y$. A picture may always be
shaded black and white with the $y$-axis being part of the boundary
of a white region.

A picture $\Theta$ will be said to be {\it standard} \ if
\begin{verse}
(i) All its cusps are in standard form, and the region
immediately below the cusp is white.

(ii) The $y$-coordinate, restricted to strings, has only
generic singularities, i.e., isolated maxima and minima.

(iii) The $y$-coordinates of all cusps and  all maxima and
minima are distinct. This set will be written
$\frak S(\Theta)=\{ y_1,y_2,y_3,\dots ,y_c\}$ with
$y_i < y_{i+1}$ for $1\leq i < c$.
\end{verse}

\ni An example of a standard $k$-picture with $k=6$ is in Figure
4.2.4.
%%%%%%%%%%%%%%%%%%%%%%%%%%%%%%%%%%%%%%%%%%%%%%%%%%%%%%%%%%%%%%%%%%%%%%%%%%%
\[
	\begin{picture}(0,0)%
\epsfig{file=ps/pic86.ps}%
\end{picture}%
\setlength{\unitlength}{0.00041700in}%
%
\begingroup\makeatletter\ifx\SetFigFont\undefined
% extract first six characters in \fmtname
\def\x#1#2#3#4#5#6#7\relax{\def\x{#1#2#3#4#5#6}}%
\expandafter\x\fmtname xxxxxx\relax \def\y{splain}%
\ifx\x\y   % LaTeX or SliTeX?
\gdef\SetFigFont#1#2#3{%
  \ifnum #1<17\tiny\else \ifnum #1<20\small\else
  \ifnum #1<24\normalsize\else \ifnum #1<29\large\else
  \ifnum #1<34\Large\else \ifnum #1<41\LARGE\else
     \huge\fi\fi\fi\fi\fi\fi
  \csname #3\endcsname}%
\else
\gdef\SetFigFont#1#2#3{\begingroup
  \count@#1\relax \ifnum 25<\count@\count@25\fi
  \def\x{\endgroup\@setsize\SetFigFont{#2pt}}%
  \expandafter\x
    \csname \romannumeral\the\count@ pt\expandafter\endcsname
    \csname @\romannumeral\the\count@ pt\endcsname
  \csname #3\endcsname}%
\fi
\fi\endgroup
\begin{picture}(7072,4244)(151,-4283)
\put(676,-3736){\makebox(0,0)[lb]{\smash{\SetFigFont{12}{14.4}{rm}$y_1$}}}
\put(676,-3136){\makebox(0,0)[lb]{\smash{\SetFigFont{12}{14.4}{rm}$y_2$}}}
\put(676,-2761){\makebox(0,0)[lb]{\smash{\SetFigFont{12}{14.4}{rm}$y_3$}}}
\put(151,-2236){\makebox(0,0)[lb]{\smash{\SetFigFont{12}{14.4}{rm}$y_5$}}}
\put(676,-811){\makebox(0,0)[lb]{\smash{\SetFigFont{12}{14.4}{rm}$y_9$}}}
\put(676,-1561){\makebox(0,0)[lb]{\smash{\SetFigFont{12}{14.4}{rm}$y_7$}}}
\put(676,-2386){\makebox(0,0)[lb]{\smash{\SetFigFont{12}{14.4}{rm}$y_4$}}}
\put(676,-1186){\makebox(0,0)[lb]{\smash{\SetFigFont{12}{14.4}{rm}$y_8$}}}
\put(676,-1936){\makebox(0,0)[lb]{\smash{\SetFigFont{12}{14.4}{rm}$y_6$}}}
\put(526,-436){\makebox(0,0)[lb]{\smash{\SetFigFont{12}{14.4}{rm}$y_{10}$}}}
\end{picture}

\]
\begin{center}
	Figure 4.2.4:  A standard 6-picture with
$\frak S(\Theta)=\{ y_1,y_2,y_3,\dots ,y_{10}\}$.
\end{center}
%%%%%%%%%%%%%%%%%%%%%%%%%%%%%%%%%%%%%%%%%%%%%%%%%%%%%%%%%%%%%%%
\bs
A standard $k$ picture $\Theta$ will be {\it labelled} \ if there is a
function from the cusps of $\Theta$ to $\coprod_m V_m$ so that a
cusp where $2m$ strings meet is assigned an element of
$V_m$ (or $N'\cap M_{m-1}$, via $\th$). We now describe  how to
associate an element $Z_{\Theta}$ of $N'\cap M_{k-1}$ to a
labelled standard $k$-picture $\Theta$, using the operators of
definition 4.1.19.

Let $\Theta$ be a labelled standard $k$-picture with
$\frak S(\Theta)=\{ y_1,y_2,y_3,\dots ,y_c\}$. We define a locally
constant function $Z:[0,1]\backslash\frak S(\Theta)\to V_{k(y)}$,
where $k(y)$ is the number  of distinct intervals of
$[0,1]\times \{y\}$ which are the connected components of its
intersection with all the black regions of $\Theta$. For instance, if
$\Theta$ is as in Figure 4.2.4 and $y_8 < y < y_9$ then $k(y)=4$.
Obviously $k(y)$ is locally constant and $k(y)=k$ for $y_c < y < 1$.
For $y < y_1$, \ $V_0=\Bbb C$ and we set $Z(y)=1$. There are five
possible ways for $Z$ to change as $y$ goes from a value $y_-$,
just less than $y_i$, to $y_+$, just bigger than $y_i$. We define
$Z(y_+)$ from $Z(y_-)$ in each case:

\medskip${\underline{\mbox{Case (i)}}}$. $y_i$ is the
$y$-coordinate of a cusp, between the $(j-1)^{\text{th}}$ and
$j^{\text{th}}$ connected components of the intersection of
$[0,1]\x y_-$ with the black regions, as below
\[
	\begin{picture}(0,0)%
\epsfig{file=ps/pic87.ps}%
\end{picture}%
\setlength{\unitlength}{0.00041700in}%
%
\begingroup\makeatletter\ifx\SetFigFont\undefined
% extract first six characters in \fmtname
\def\x#1#2#3#4#5#6#7\relax{\def\x{#1#2#3#4#5#6}}%
\expandafter\x\fmtname xxxxxx\relax \def\y{splain}%
\ifx\x\y   % LaTeX or SliTeX?
\gdef\SetFigFont#1#2#3{%
  \ifnum #1<17\tiny\else \ifnum #1<20\small\else
  \ifnum #1<24\normalsize\else \ifnum #1<29\large\else
  \ifnum #1<34\Large\else \ifnum #1<41\LARGE\else
     \huge\fi\fi\fi\fi\fi\fi
  \csname #3\endcsname}%
\else
\gdef\SetFigFont#1#2#3{\begingroup
  \count@#1\relax \ifnum 25<\count@\count@25\fi
  \def\x{\endgroup\@setsize\SetFigFont{#2pt}}%
  \expandafter\x
    \csname \romannumeral\the\count@ pt\expandafter\endcsname
    \csname @\romannumeral\the\count@ pt\endcsname
  \csname #3\endcsname}%
\fi
\fi\endgroup
\begin{picture}(6623,3543)(901,-3883)
\put(901,-1036){\makebox(0,0)[lb]{\smash{\SetFigFont{12}{14.4}{rm}$y_+$}}}
\put(901,-2086){\makebox(0,0)[lb]{\smash{\SetFigFont{12}{14.4}{rm}$y_i$}}}
\put(901,-3136){\makebox(0,0)[lb]{\smash{\SetFigFont{12}{14.4}{rm}$y_-$}}}
\put(3151,-3811){\makebox(0,0)[lb]{\smash{\SetFigFont{12}{14.4}{rm}$j-1$}}}
\put(5701,-3811){\makebox(0,0)[lb]{\smash{\SetFigFont{12}{14.4}{rm}$j$}}}
\end{picture}

\]

\bs

\ni Set \ $Z(y_+)=\a_{j,c}(Z(y_-))$ where $c$ is the label
(in $V_{k(y_-)}$) associated with the cusp.

\bs ${\underline{\mbox{Case (ii)}}}$. $y_i$ is the
$y$-coordinate of a minimum, with numbering and shading as below:
\[
	\begin{picture}(0,0)%
\epsfig{file=ps/pic88.ps}%
\end{picture}%
\setlength{\unitlength}{0.00041700in}%
%
\begingroup\makeatletter\ifx\SetFigFont\undefined
% extract first six characters in \fmtname
\def\x#1#2#3#4#5#6#7\relax{\def\x{#1#2#3#4#5#6}}%
\expandafter\x\fmtname xxxxxx\relax \def\y{splain}%
\ifx\x\y   % LaTeX or SliTeX?
\gdef\SetFigFont#1#2#3{%
  \ifnum #1<17\tiny\else \ifnum #1<20\small\else
  \ifnum #1<24\normalsize\else \ifnum #1<29\large\else
  \ifnum #1<34\Large\else \ifnum #1<41\LARGE\else
     \huge\fi\fi\fi\fi\fi\fi
  \csname #3\endcsname}%
\else
\gdef\SetFigFont#1#2#3{\begingroup
  \count@#1\relax \ifnum 25<\count@\count@25\fi
  \def\x{\endgroup\@setsize\SetFigFont{#2pt}}%
  \expandafter\x
    \csname \romannumeral\the\count@ pt\expandafter\endcsname
    \csname @\romannumeral\the\count@ pt\endcsname
  \csname #3\endcsname}%
\fi
\fi\endgroup
\begin{picture}(6623,3544)(901,-3883)
\put(901,-1036){\makebox(0,0)[lb]{\smash{\SetFigFont{12}{14.4}{rm}$y_+$}}}
\put(901,-2086){\makebox(0,0)[lb]{\smash{\SetFigFont{12}{14.4}{rm}$y_i$}}}
\put(901,-3136){\makebox(0,0)[lb]{\smash{\SetFigFont{12}{14.4}{rm}$y_-$}}}
\put(3151,-3811){\makebox(0,0)[lb]{\smash{\SetFigFont{12}{14.4}{rm}$j-1$}}}
\put(4501,-3811){\makebox(0,0)[lb]{\smash{\SetFigFont{12}{14.4}{rm}$j$}}}
\put(5701,-3811){\makebox(0,0)[lb]{\smash{\SetFigFont{12}{14.4}{rm}$j+1$}}}
\end{picture}

\]
\bs

\ni Set \ $Z(y_+)=\kappa_j(Z(y_-))$.

\newpage ${\underline{\mbox{Case (iii)}}}$. $y_i$ is the
$y$-coordinate of a minimum, with numbering and shading as below:
\[
	\begin{picture}(0,0)%
\epsfig{file=ps/pic89.ps}%
\end{picture}%
\setlength{\unitlength}{0.00041700in}%
%
\begingroup\makeatletter\ifx\SetFigFont\undefined
% extract first six characters in \fmtname
\def\x#1#2#3#4#5#6#7\relax{\def\x{#1#2#3#4#5#6}}%
\expandafter\x\fmtname xxxxxx\relax \def\y{splain}%
\ifx\x\y   % LaTeX or SliTeX?
\gdef\SetFigFont#1#2#3{%
  \ifnum #1<17\tiny\else \ifnum #1<20\small\else
  \ifnum #1<24\normalsize\else \ifnum #1<29\large\else
  \ifnum #1<34\Large\else \ifnum #1<41\LARGE\else
     \huge\fi\fi\fi\fi\fi\fi
  \csname #3\endcsname}%
\else
\gdef\SetFigFont#1#2#3{\begingroup
  \count@#1\relax \ifnum 25<\count@\count@25\fi
  \def\x{\endgroup\@setsize\SetFigFont{#2pt}}%
  \expandafter\x
    \csname \romannumeral\the\count@ pt\expandafter\endcsname
    \csname @\romannumeral\the\count@ pt\endcsname
  \csname #3\endcsname}%
\fi
\fi\endgroup
\begin{picture}(6623,3543)(901,-3883)
\put(901,-1036){\makebox(0,0)[lb]{\smash{\SetFigFont{12}{14.4}{rm}$y_+$}}}
\put(901,-2086){\makebox(0,0)[lb]{\smash{\SetFigFont{12}{14.4}{rm}$y_i$}}}
\put(901,-3136){\makebox(0,0)[lb]{\smash{\SetFigFont{12}{14.4}{rm}$y_-$}}}
\put(3151,-3811){\makebox(0,0)[lb]{\smash{\SetFigFont{12}{14.4}{rm}$j-1$}}}
\put(5776,-3811){\makebox(0,0)[lb]{\smash{\SetFigFont{12}{14.4}{rm}$j$}}}
\end{picture}

\]

\bs

\ni Set \ $Z(y_+)=\eta_j(Z(y_-))$.


\bs ${\underline{\mbox{Case (iv)}}}$. $y_i$ is the
$y$-coordinate of a maximum, with numbering and shading as below:
\[
	\begin{picture}(0,0)%
\epsfig{file=ps/pic90.ps}%
\end{picture}%
\setlength{\unitlength}{0.00041700in}%
%
\begingroup\makeatletter\ifx\SetFigFont\undefined
% extract first six characters in \fmtname
\def\x#1#2#3#4#5#6#7\relax{\def\x{#1#2#3#4#5#6}}%
\expandafter\x\fmtname xxxxxx\relax \def\y{splain}%
\ifx\x\y   % LaTeX or SliTeX?
\gdef\SetFigFont#1#2#3{%
  \ifnum #1<17\tiny\else \ifnum #1<20\small\else
  \ifnum #1<24\normalsize\else \ifnum #1<29\large\else
  \ifnum #1<34\Large\else \ifnum #1<41\LARGE\else
     \huge\fi\fi\fi\fi\fi\fi
  \csname #3\endcsname}%
\else
\gdef\SetFigFont#1#2#3{\begingroup
  \count@#1\relax \ifnum 25<\count@\count@25\fi
  \def\x{\endgroup\@setsize\SetFigFont{#2pt}}%
  \expandafter\x
    \csname \romannumeral\the\count@ pt\expandafter\endcsname
    \csname @\romannumeral\the\count@ pt\endcsname
  \csname #3\endcsname}%
\fi
\fi\endgroup
\begin{picture}(6623,3544)(901,-3883)
\put(901,-1036){\makebox(0,0)[lb]{\smash{\SetFigFont{12}{14.4}{rm}$y_+$}}}
\put(901,-2086){\makebox(0,0)[lb]{\smash{\SetFigFont{12}{14.4}{rm}$y_i$}}}
\put(901,-3136){\makebox(0,0)[lb]{\smash{\SetFigFont{12}{14.4}{rm}$y_-$}}}
\put(5776,-3811){\makebox(0,0)[lb]{\smash{\SetFigFont{12}{14.4}{rm}$j+1$}}}
\put(3226,-3811){\makebox(0,0)[lb]{\smash{\SetFigFont{12}{14.4}{rm}$j$}}}
\end{picture}

\]

\bs

\ni Set \ $Z(y_+)=\mu_{j+1}(Z(y_-))$.

\bs ${\underline{\mbox{Case (v)}}}$. $y_i$ is the
$y$-coordinate of a maximum  with shading as below:
\[
	\begin{picture}(0,0)%
\epsfig{file=ps/pic91.ps}%
\end{picture}%
\setlength{\unitlength}{0.00041700in}%
%
\begingroup\makeatletter\ifx\SetFigFont\undefined
% extract first six characters in \fmtname
\def\x#1#2#3#4#5#6#7\relax{\def\x{#1#2#3#4#5#6}}%
\expandafter\x\fmtname xxxxxx\relax \def\y{splain}%
\ifx\x\y   % LaTeX or SliTeX?
\gdef\SetFigFont#1#2#3{%
  \ifnum #1<17\tiny\else \ifnum #1<20\small\else
  \ifnum #1<24\normalsize\else \ifnum #1<29\large\else
  \ifnum #1<34\Large\else \ifnum #1<41\LARGE\else
     \huge\fi\fi\fi\fi\fi\fi
  \csname #3\endcsname}%
\else
\gdef\SetFigFont#1#2#3{\begingroup
  \count@#1\relax \ifnum 25<\count@\count@25\fi
  \def\x{\endgroup\@setsize\SetFigFont{#2pt}}%
  \expandafter\x
    \csname \romannumeral\the\count@ pt\expandafter\endcsname
    \csname @\romannumeral\the\count@ pt\endcsname
  \csname #3\endcsname}%
\fi
\fi\endgroup
\begin{picture}(6623,3544)(901,-3883)
\put(901,-1036){\makebox(0,0)[lb]{\smash{\SetFigFont{12}{14.4}{rm}$y_+$}}}
\put(901,-2086){\makebox(0,0)[lb]{\smash{\SetFigFont{12}{14.4}{rm}$y_i$}}}
\put(901,-3136){\makebox(0,0)[lb]{\smash{\SetFigFont{12}{14.4}{rm}$y_-$}}}
\put(2626,-3811){\makebox(0,0)[lb]{\smash{\SetFigFont{12}{14.4}{rm}$j-1$}}}
\put(4351,-3811){\makebox(0,0)[lb]{\smash{\SetFigFont{12}{14.4}{rm}$j$}}}
\put(6226,-3811){\makebox(0,0)[lb]{\smash{\SetFigFont{12}{14.4}{rm}$j+1$}}}
\end{picture}

\]

\bs

\ni Set \ $Z(y_+)=a_j(Z(y_-))$.

\medskip
Finally, we define \ $Z_{\Theta}=\th(Z(y))$ for $y > y_c$,
also written $Z_1(\Theta)$.
Our main job is now to prove that $Z_{\Theta}$ is unchanged if $\Theta$
is changed by isotopy to another standard picture $\Theta'$, with
labels transported by the isotopy. If the isotopy passes only
through standard pictures, critical points can never change order
or be annihilated or created. The pattern of connected components of
the intersection of the black regions with horizontal lines cannot
be changed either as such a change would have to involve two
maxima, minima or cusps having the same $y$ coordinate. So
isotopies through standard pictures do not change $Z_{\Theta}$.

We next argue that if $\phi_t$, $0\leq t\leq 1$ is an isotopy which
preserves the standard form of each cusp, then
$Z_{\Theta}=Z_{\phi_1(\Theta)}$. For now $Z_{\phi_t(\Theta)}$ can
only change if the singularities of the $y$-coordinate function
change. By putting the isotopy in general position we see that this
can be supposed to happen in only two ways (see e.g. [Tu], or
note that this argument can be made quite combinatorial by using
piecewise linear strings).



